\section*{TLDR; Tjekliste til Bandboss}

\textbf{Sådan laver du din første revy som bandboss}

\subsection*{Inden revyperioden}

\begin{enumerate}
\item \textbf{Vælg materiale}\\
Afhængig af hvordan din revy er struktureret, har du forskellig indflydelse på hvilke sange der skal spilles til revyen. Uanset hvordan det bliver bestemt, så hold gerne øje med at sangene bliver valgt i nogenlunde god tid, så der er tid til at forberede dem.

\item \textbf{Sæt et band}\\
Dette kan gøres på mange måder. Hold auditions for hvert instrument, eller hold en jam-dag hvor alle interesserede kan møde op og spille sammen, så du kan få et indtryk af hvordan folk spiller sammen. Dette skal gerne gøres 6-8 uger inden revyen.\\
Rækkefølgen på trin 1 og 2 kan godt byttes om.

\item \textbf{Bestem tonearter på sange}\\
Hvis din revy har en sangboss, tager denne formentlig ansvaret for at sætte toneleje for sangene ud fra foresyngning med sangerne. Pro tips for optimale valg af tonearter findes i kapitlet derom.

\item \textbf{Fordel ansvar for arrangementer}\\
Til en revy indenfor SaTyR-rammen er der typisk 14-16 sange i revyen. Det er mere end en person kan skrive arrangementer til, i hvert fald hvis man også skal læse og sove. Fordel ansvaret så resten af bandet også hjælper. Ofte er rytmegruppe-medlemmer ferme til dette, men blæsere kan og burde også aktiveres.

\item \textbf{Skriv arrangementer}\\
Det er ikke livsnødvendigt at disse er klar når man begynder øvning, men det gør alle numre bedre at man har en vision og et mål med sin øvning. Nu tidligere det bliver klart, desto mere tid har man til at realisere det.

\item \textbf{Afhold indledende øvere}\\
Inden den skemalagte revyøvning for alle revyster starter, er det godt at have øvet sangene med bandet. På den måde får man mest muligt ud af den dyrebare øvetid under revyen. Afhold ca. 3 øvere inden revyperioden starter, og prøv at få spillet ca 5 sange til hver. På den måde sørger man for at bandet kan formen og akkorderne til sangene inden øvning med sangere. Så er der god tid til at arbejde med lir efterfølgende.

\item \textbf{Kommuniker med TeXnikken} - hviken besætning har i og hvilke outputs/mics skal i bruge?

\item \textbf{Find ekstern instruktør} til at se kunstnerisk gennemspilning.\\
De fleste revyer har en gennemspilning hvor et panel af erfarne revyster giver feedback til sketches og sange, for at få et blik udefra. Dette kan også være meget gavnligt for bandet, så overvej at finde en der vil give feedback til bandet specifikt.\\
Disse har ofte også feedback til lyden, som er svær at høre fra bandscenen, og det kan være meget værdifuldt.

\item \textbf{Lav noget hyggeligt som band} - tag på cafeen efter øver, spis morgenmad sammen i starten af en øveweekend eller noget andet. Godt sammenhold giver god musik (:
\end{enumerate}

\subsection*{Under revyperioden}

\begin{enumerate}
\item \textbf{Kommuniker med, og hold styr på bandet}\\
Er der nogen der kun kan være der til nogle af dagene/noget af tiden? Sørg i så fald for at øvning bliver planlagt så musikerne er til stede på de numre de er vigtige på. Sørg også for at bandet ved hvor og hvornår de skal møde, og hvor længe de forventes at blive.

\item \textbf{Giv input til dem, der lægger øveplan}\\
Bed om flere øvesegmenter til sange der er svære af den ene eller den anden grund.

\item \textbf{Sørg for at bandet får pauser og snacks}

\item \textbf{Bestem cues til sange} - hvordan skal sangene starte, og hvordan ved man at alle er klar? (Mere om dette i sit eget kapitel)

\item \textbf{Vælg bandtema} og skriv/arranger bandintroer.

\item \textbf{Bed om at bandet kan spise først}, da de typisk har stramt program både før og efter aftensmad
\end{enumerate}

\subsection*{Under forestillingerne}

\begin{enumerate}
\item Sørg for at bandet er klar til at gå på når det er showtime.
Det er altid bandet der starter, så det er MEGET vigtigt at alle bandmedlemmer er klar et par minutter inden. Dette gælder inden alle akter. Sørg også for at kommunikere med TeXnikken/koordinatorer om hvornår man starter efter pauser, og om der er forsinkelser.

\item Tag cues fra scenen og giv signaler til at starte sange.
Dette kan evt. uddelegeres afhængig af opsætning på bandscenen.
\end{enumerate}

\subsection*{Til efterfest}

\begin{enumerate}
\item Hav det sjovt

\item Lad være med at brække arme/fingre (i hvert fald inden lørdag)
\end{enumerate}

\subsection*{Til oprydning}

\begin{enumerate}
\item Ryd bandscenen ASAP

\item Ryd. Bandscenen. ASAP.

\item Gør dyrestalden rent

\item Gør. Dyrestalden. Rent.

\end{enumerate}

\subsection*{Generelt: Uddelegér!}
Det er ikke dig, der skal gøre det hele.
Du skal bare sørge for at hver ting bliver gjort af \textit{nogen}.